\section{Related Work}
\paragraph{\textbf{Code Generation Benchmarks.}} The evaluation of code generation has undergone a significant evolution from isolated code snippets to complex development environments. Early benchmarks, such as HumanEval~\cite{humaneval} and MBPP~\cite{austin2021program}, primarily defined evaluation as text-to-code translation tasks. These benchmarks typically provided a brief natural-language description, such as a docstring or a function signature, requiring the model to synthesize an independent function body. Subsequently, benchmarks like ClassEval~\cite{classeval} and CoderEval~\cite{yu2024codereval} introduced class-level and dependency-aware tasks. With the expansion of LLM context windows, evaluation has further advanced to the repository level. Benchmarks such as DevEval~\cite{li2024deveval} and EvoCodeBench~\cite{li2024evocodebench} focus on repository-level code generation to assess the utilization of global repository context. However, these benchmarks typically presuppose that the modification location is known by providing specific function signatures as inputs. This design fails to evaluate the agent's ability to navigate the repository autonomously. Conversely, benchmarks represented by SWE-bench~\cite{swebench} target issue resolution, requiring agents to autonomously locate and fix bugs based on issue descriptions or tracebacks. Although it includes approximately 18-22\% feature implementation scenarios~\cite{rashid2025swe}, it lacks a dedicated focus on the generative process of implementing new features from scratch. More recently, benchmarks like FEA-Bench~\cite{li2025fea} and NoCodeBench~\cite{deng2025nocode} have been proposed to specifically evaluate feature implementation. However, they rely on inputs that simplify the problem: FEA-Bench explicitly provides the signatures and documentation of the new components, and NoCodeBench utilizes documentation updates as specifications. These setups significantly reduce the difficulty of the task, thereby failing to sufficiently evaluate the performance of SOTA coding agents. While FeatBench shares the objective of evaluating feature implementation, it distinguishes itself by strictly adhering to realistic constraints without code hints. This rigorous formulation requires agents to autonomously identify appropriate injection points and implement features solely from natural language requirements.

\paragraph{\textbf{Base Models and Software Engineering Agents.}} The advancement of software engineering agents is underpinned by the rapid evolution of base LLMs. Proprietary models like GPT-5~\cite{gpt-5} have set performance milestones. Parallelly, open-source models including CodeLLaMa~\cite{roziere2023code}, Qwen3-Coder-Next~\cite{qwen3lmQwen3Coder}, and Kimi-2~\cite{team2025kimi} have demonstrated impressive proficiency by leveraging massive code corpora and specialized architectures. Building upon these powerful backbones, the architecture of agents has evolved from single generative models into complex systems. Current architectures are primarily categorized into two paradigms. Pipeline-based agents, such as Agentless~\cite{agentless}, adopt a fixed localization-repair workflow that emphasizes cost-effectiveness and stability, yet they lack the flexibility to handle implicit cross-file dependencies. In contrast, autonomous agents, such as Trae-agent~\cite{traeagent} and SWE-agent~\cite{yang2024swe}, employ a ReAct loop~\cite{yao2022react} to dynamically plan paths and utilize tools, thereby possessing stronger environment exploration capabilities. However, the existing literature comparing these two types of agents primarily focuses on bug-fixing rates. In feature implementation tasks, the unique behavioral patterns of autonomous agents, particularly the phenomena of ``scope creep'', have not been sufficiently studied.